\section{Dataset Details}
\label{app:datasets}

\subsection{FEMTO (PRONOSTIA)}

The FEMTO bearing dataset~\citep{nectoux2012pronostia} was collected at the FEMTO-ST institute using an accelerated degradation test platform. Bearings are subjected to radial loads while rotating at controlled speeds, with vibration measured by a piezoelectric accelerometer at 25.6\,kHz. The dataset contains 17 run-to-failure experiments across 3 operating conditions (we use 16 with complete recordings).

\subsection{XJTU-SY}

The XJTU-SY bearing dataset~\citep{wang2018xjtusybearing} contains 15 rolling element bearings tested under 3 operating conditions on the XJTU-SY platform. Each bearing is monitored with two accelerometers (horizontal and vertical) at 25.6\,kHz. We use 7 bearings with clear degradation trajectories.

\subsection{Pretraining Data Sources}

JEPA pretraining uses 33,939 vibration windows from 8 sources:

\begin{table}[h]
\centering
\caption{Pretraining data sources.}
\label{tab:pretrain_sources}
\small
\begin{tabular}{@{}lccl@{}}
\toprule
Source & Windows & Sampling Rate & Fault Types \\
\midrule
CWRU & 2,400 & 12\,kHz & Healthy, IR, OR, Ball \\
MFPT & 1,200 & 25.6\,kHz & Healthy, IR, OR \\
IMS & 3,000 & 20\,kHz & Run-to-failure \\
FEMTO & 8,000 & 25.6\,kHz & Run-to-failure \\
XJTU-SY & 4,500 & 25.6\,kHz & Run-to-failure \\
Paderborn & 6,000 & 64\,kHz & Healthy, IR, OR \\
Ottawa & 4,839 & 12\,kHz & Healthy, IR, OR, Ball \\
MAFAULDA & 4,000 & 50\,kHz & Multiple \\
\bottomrule
\end{tabular}
\end{table}

All data is resampled to a common rate before pretraining.

\section{Hyperparameters}
\label{app:hyperparams}

\begin{table}[h]
\centering
\caption{Complete hyperparameter settings for all methods.}
\label{tab:hyperparams}
\small
\begin{tabular}{@{}llc@{}}
\toprule
Component & Hyperparameter & Value \\
\midrule
\multirow{7}{*}{JEPA Pretraining} & Encoder dim & 256 \\
& Encoder depth & 4 layers \\
& Attention heads & 4 \\
& Predictor dim & 128 \\
& Predictor depth & 4 layers \\
& Mask ratio & 62.5\% \\
& EMA momentum & 0.996 \\
& Learning rate & $1 \times 10^{-4}$ \\
& Batch size & 256 \\
& Epochs & 100 (best at epoch 2) \\
& Loss & L1 + variance reg ($\lambda = 0.1$) \\
\midrule
\multirow{4}{*}{Contrastive} & Learning rate & $1 \times 10^{-4}$ \\
& Margin $\alpha$ & 1.0 \\
& Epochs & 100 \\
& Episodes & 18 (training split) \\
\midrule
\multirow{4}{*}{RUL Head (LSTM)} & Hidden dim & 128 \\
& Layers & 2 \\
& Dropout & 0.2 \\
& Learning rate & $1 \times 10^{-3}$ \\
\midrule
\multirow{4}{*}{RUL Head (Transformer)} & Layers & 2 \\
& Heads & 4 \\
& Dim & 256 \\
& Learning rate & $1 \times 10^{-3}$ \\
\bottomrule
\end{tabular}
\end{table}

\section{Additional Results}
\label{app:additional}

\subsection{Piecewise Label Ablation}

\begin{table}[h]
\centering
\caption{Linear vs.\ piecewise-linear RUL labels. Piecewise labels define a constant healthy phase (RUL=1) followed by linear degradation, making prediction harder due to ambiguity in the healthy phase.}
\label{tab:piecewise}
\small
\begin{tabular}{@{}lccc@{}}
\toprule
Method & Linear RMSE & Piecewise RMSE & $\Delta$ \\
\midrule
JEPA + LSTM & 0.189 & 0.310 & +0.121 \\
HC + MLP & 0.085 & 0.114 & +0.029 \\
HC + LSTM & 0.177 & 0.214 & +0.037 \\
Transformer + HC & 0.070 & 0.114 & +0.044 \\
\bottomrule
\end{tabular}
\end{table}

\subsection{JEPA Pretraining Configurations}

\begin{table}[h]
\centering
\caption{Pretraining configurations explored. Best validation loss does not always correspond to best downstream RUL performance.}
\label{tab:configs}
\small
\begin{tabular}{@{}lccccl@{}}
\toprule
Config & LR & EMA & Mask & Best Val Loss & Verdict \\
\midrule
Default & $10^{-4}$ & 0.996 & 62.5\% & 0.0166 (ep.\ 2) & \textbf{Best} \\
Lower LR & $3 \times 10^{-5}$ & 0.996 & 62.5\% & 0.0172 & Similar \\
Fixed LR & $10^{-5}$ & 0.996 & 62.5\% & 0.0189 & Worse \\
High EMA & $10^{-4}$ & 0.999 & 62.5\% & 0.0178 & Worse \\
2ch FFT & $10^{-4}$ & 0.996 & 62.5\% & 0.0149 & Better loss, worse RUL \\
\bottomrule
\end{tabular}
\end{table}
